\begin{vidu}%[3P7-2.2B][Loai1]
Cho khối lượng của hạt nhân $\Hn{84}{210}{Po}$ là $209,9828$u; của neutron là 1,0087u; của proton là 1,0073u. Cho $1\ uc^2=931,5\ MeV$ Xét hạt nhân \Hn{84}{210}{Po}
\begin{itemize}
\item[a)] Tính độ hụt khối của hạt nhân.
\item[b)] Tính năng lượng liên kết, năng lượng liên kết riêng. 
\item[c)] Các hạt nhân deuterium, \Hn{2}{4}{He}, \Hn{53}{139}{I}, \Hn{92}{235}{U}
có khối lượng tương ứng là 2,0136u, 4,0015u; 138,8970u và 234,9933u. Sắp xếp các hạt nhân trên theo thứ tự giảm dần về độ bền vững của hạt nhân.
\end{itemize}
\loigiai{
\begin{itemize}
\item[a)] Độ hụt khối của hạt nhân \Hn{84}{210}{Po}: $\Delta m = 84m_p + (210 - 84)m_n - m_{Po} = 1,7266$u.
\item[b)] 
\begin{itemize}
\item Năng lượng liên kết:
\begin{align*}
W_{lk}&=\left(84.m_p+(210-84)m_n-m_{Po}\right)c^2=\left(84.1,0073+(210-84).1,0087-209,9828\right)uc^2\\
&=1,7266.931,5\ MeV=1608,3\ MeV.
\end{align*}
\item Năng lượng liên kết riêng: $W_{lkr}=\dfrac{W_{lk}}{A}=\dfrac{1608,3}{210}=7,66\ MeV/\text{nuclôn}$.
\end{itemize}	
\item[c)] 
\begin{itemize}
\item $W_{lkrHe}=\dfrac{W_{lkHe}}{A_{He}}=1,1\ MeV$.
\item $\Delta m_{He}=0,0305u\Rightarrow W_{lkHe}=\Delta m_{He}. c^2=28,410\Rightarrow W_{lkrHe}=\dfrac{W_{lkHe}}{A_{He}}=\dfrac{28,410}{4}=7,1\ MeV$.
\item $\Delta {m_I}=1,2381u\Rightarrow W_{lkI}=\Delta {m_I}. c^2=1153,29\Rightarrow W_{lkrI}=\dfrac{W_{lkI}}{A_I}=\dfrac{1153,29}{139}=8,29\ MeV$.
\item $\Delta {m_U}=1,9224u\Rightarrow W_{lkU}=\Delta {m_U}. c^2=1790,71\Rightarrow W_{lkrU}=\dfrac{W_{lkU}}{A_U}=\dfrac{1790,71}{235}=7,6\ MeV$.
\item Thứ tự giảm dần về độ bền vững là: I, U, He.
\end{itemize}
\end{itemize}
}
\end{vidu}
\loai{Tính độ hụt khối}
\begin{vidu}%[3P7-2.2B][Loai1] 
\nguon{QG - 2019} Cho khối lượng của proton, neutron, hạt nhân $_3^6Li$ lần lượt là 1,0073u; 1,0087u; 6,0135u. Độ hụt khối của $_3^6Li$ là
\choice
{\True 0,0345u}
{0,0245u}
{0,0512u}
{0,0412u}
\loigiai{Độ hụt khối hạt nhân Li:
$\Delta m=Z.m_p+\left(A-Z\right)m_n-m_X=3.1,0073u+3.1,0087u-6,0135u=0,0345u.$}
\end{vidu}
\begin{vidu}%[3P7-2.2B][Loai1]
\nguon{QG - 2017} Hạt nhân $\Hn{8}{17}{O}$ có khối lượng 16,9947u. Biết khối lượng của proton và neutron lần lượt là 1,0073u và 1,0087u. Độ hụt khối của $\Hn{8}{17}{O}$ là	
\choice
{0,1406u}
{0,1294u}
{0,1532u}
{\True 0,1420u}
\loigiai{
Độ hụt khối của $\Hn{8}{17}{O}$
\begin{align*}
\Delta m=Z.m_p+\left(A-Z\right)m_n-m=8.1,0073+\left(17-8\right).1,0087-16,9947=0,142u.
\end{align*}
}
\end{vidu}
\loai{Tính năng lượng liên kết}
\begin{vidu} %[3P7-2.2B][Loai2]
\nguon{QG - 2020} 	Hạt nhân $_{47}^{107}Ag$ có khối lượng $106,8783u $. Cho khối lượng của protton và noton lần lượt là $1,0073u$ và $1,0087u$; $1u=931,5\ MeV/c^2$. Năng lượng liên kết của hạt nhân $_{47}^{107}Ag$ là
\choice
{\True $919,2\ MeV$}
{$902,3\ MeV$}
{$938,3\ MeV$}
{$939,6\ MeV$}
\loigiai{
Năng lượng liên kết của hạt nhân $_{47}^{107}Ag$:
\begin{align*}
W_{lk}=\left[{Z.m_p+(A-Z).m_N-m}\right].931,5=[47.1,0073+60.1,0087-106,8783].931,5=919,2042\ MeV
\end{align*}}
\end{vidu}
\loai{Tính năng lượng liên kết riêng}
\begin{vidu}%[3P7-2.2B][Loai3]
\nguon{ĐH - 2008} Hạt nhân $\Hn{4}{10}{Be}$ có khối lượng $10,0135u$. Khối lượng của neutron $m_n = 1,0087u$, khối lượng của proton $m_p = 1,0073u,\ 1u = 931\ MeV/c^2$. Năng lượng liên kết riêng của hạt nhân $\Hn{4}{10}{Be}$ là
\choice
{0,6321 MeV}
{63,2152 MeV}
{\True 6,3215 MeV}
{632,1531 MeV}
\loigiai{
Ta có: $W_{lkr}=\dfrac{W_{\ell k}}{A}=\dfrac{\left(4m_p+6m_n-m_{Be}\right)c^2}{10}=6,3215\ MeV$.
}
\end{vidu}
\loai{Cho năng lượng liên kết tìm năng lượng liên kết riêng}
\begin{vidu}%[3P7-2.2B][Loai4]
\nguon{QG - 2018} Hạt nhân $_{92}^{235}U$ có năng lượng liên kết là 1784 MeV. Năng lượng liên kết riêng của hạt nhân này là
\choice
{5,45 MeV/nuclôn}
{12,47 MeV/nuclôn}
{\True 7,59 MeV/nuclôn}
{19,39 MeV/nuclôn}
\loigiai{
Ta có: $W_{lkr}=\dfrac{W_{lk}}{A}=\dfrac{1784}{235}=7,59\ MeV/\text{nuclôn}$.
}
\end{vidu}
\loai{So sánh tính bền vững của 2 hạt nhân}
\begin{vidu}%[3P7-2.2K][Loai5]
\nguon{ĐH - 2011} Cho khối lượng của proton, neutron, $\Hn{18}{40}{Ar},\ \Hn{3}{6}{Li}$ lần lượt là: 1,0073u; 1,0087u; 39,9525u; 6,0145u và $1u = 931,5$ MeV/$c^2$. So với năng lượng liên kết riêng của hạt nhân $\Hn{3}{6}{Li}$ thì năng lượng liên kết riêng của hạt nhân $\Hn{18}{40}{Ar}$
\choice
{lớn hơn một lượng là 5,20 MeV}
{\True lớn hơn một lượng là 3,42 MeV}
{nhỏ hơn một lượng là 3,42 MeV}
{nhỏ hơn một lượng là 5,20 MeV}
\loigiai{
Tính được năng lượng liên kết riêng của Ar và Li lần lượt là 8,62 MeV và 5,20 MeV.
}
\end{vidu}
\loai{So sánh tính bền vững của nhiều hạt nhân}
\begin{vidu}%[3P7-2.2K][Loai6]
\nguon{ĐH - 2012} Các hạt nhân deuterium $_1^2H$; tritium $_1^3H$, helium $_2^4He$ có năng lượng liên kết lần lượt là 2,22 MeV; 8,49 MeV và 28,16 MeV. Các hạt nhân trên được sắp xếp theo thứ tự giảm dần về độ bền vững của hạt nhân là
\choice
{$_1^2H;\ _2^4He;\ _1^3H$}
{$_1^2H;\ _1^3H;\ _2^4He$}
{\True $_2^4He;\ _1^3H$;\ $_1^2H$}
{$_1^3H;\ _2^4He;\ _1^2H$}
\loigiai{
\begin{itemize}
\item Năng lượng liên kết riêng của các hạt nhân deuterium $_1^2H$; tritium $_1^3H$, helium $_2^4He$ lần lượt là 1,11 MeV; 2,83 MeV; 7,04 MeV.
\item Năng lượng liên kết riêng càng lớn thì hạt càng bền vững.
\end{itemize}
}
\end{vidu}
\loai{Cho năng lượng liên kết riêng - tìm khối lượng}
\begin{vidu}%[3P7-2.2K][Loai7]
Một hạt nhân có 8 proton và 9 neutron, năng lượng liên kết riêng của hạt nhân này là 7,75 MeV/nuclôn. Biết $m_p = 1,0073$u; $m_n = 1,0087$u, $1u=931,5\ MeV/c^2$. Khối lượng của hạt nhân đó là
\choice
{16,9455u}
{17,0053u}
{\True 16,9953u}
{17,0567u}
\loigiai{
Ta có: 
\begin{align*}
W_{lkr}.A &= \left(Z.m_p + N.m_n - m\right).c^2\\
\Leftrightarrow 8,32.10^{-3}uc^2.(8 + 9) &= (8.1,0073u + 9.1,0087u - m).c^2\Rightarrow m = 16,9953u.
\end{align*}
}
\end{vidu}
\loai{Tổng hợp}
\begin{vidu}%[3P7-2.2K][Loai8]
\nguon{MH - 2017} Cho khối lượng của hạt nhân $_2^4He$; proton và neutron lần lượt là 4,0015u; 1,0073u và 1,0087u. Lấy $1u = 1,66.10^{-27}\ kg$; $c = 3.10^8\ m/s$; $N_A = 6,02.10^{23}\ mol^{-1}$. Năng lượng tỏa ra khi tạo thành 1 mol $_2^4He$ từ các nuclôn là
\choice
{$2,74.10^6$ J}
{\True $2,74.10^{12}$ J}
{$1,71.10^6$ J}
{$1,71.10^{12}$ J}
\loigiai{
\begin{itemize}
\item Ta có: $2_1^1p+2_0^1n\rightarrow_2^4He$
\item $W_{\ell k}=\left(2m_p+2m_n-m_{He}\right)c^2=28,41075\ MeV$
\item 1 mol $Fe\Rightarrow W=N_A.\Delta W= 1,71.10^{25}\ MeV=2,736.10^{12}\ J$.
\end{itemize}
}
\end{vidu}
